% New native author theorem, 2026-10-21.  Scratch draft for independent reconstruction.
\section{A physical boundary stopping proof of the full word-pressure bound}
\label{ps:section}

The capacity moment in Section~\ref{tc:section} does not, by itself,
pay the small-capacity part of the original physical-word pressure.
We now bound that physical-word sum by a different geometric
mechanism, choosing one nearly maximizing past-strip class per
physical word.  We do not assert a uniform sum over every
homogeneous class.  We stop an actual affine
transversal inside the same complete physical word and the same
past homogeneity labels, and transport its word weight to a genuine
one-dimensional source at a small positive interior level.  The
possible exits are exhaustively classified below.  No lower bound
on a cell's area, length, capacity, or angle is required.

All collisions, their lifted centers, the radius
$R\in[R_-,R_+]=[9/20,47/100]$, and the physical unit-speed flow are
unchanged.  The stable cocycle is Euclidean arclength in $(s,\varphi)$,
where $s=R\alpha$.  Put
\[
 \kappa_R=R^{-1},\qquad
 \Gamma=\sqrt{1+(20/9)^2},\qquad
 \sigma=47/59.
\]
Use the original $q=2$ homogeneity strips and their fixed $k_0$.
Their interior boundary levels form the set
\[
 \mathcal H=\{\,\pm(\pi/2-k^{-2}):k\ge k_0\,\},\qquad
 S_{\mathcal H}:=\sum_{h\in\mathcal H}\cos h
     =2\sum_{k\ge k_0}\sin(k^{-2})
     \le {2\over k_0-1}.
\]
This is a sum of actual endpoint incidence factors, not a sum of
strip widths.

We use the horizontal collision-source constants of
Theorem~\ref{thm:cs-half-sup}, writing
$H_{\rm H}=H_{\rm hor}$ and $E_{\rm H}=E_{\rm cs}=47/6$: for every prescribed
$\varphi_i=h\in(-\pi/2,\pi/2)$ and either time direction, the
source parameter is $\alpha_i$ on its one physical circle.
For $l\ge1$ the genuine branch partition has inverse terminal
momentum-derivative weights $a_I$ with
\begin{equation}\label{ps:horizontal-input}
 \sum_Ia_I\le H_{\rm H},\qquad
 a_I\le E_{\rm H}\sigma^{l-1}.
\end{equation}
The initial source need not be homogeneous.  There is no terminal
homogeneity subdivision.  These are the unconditional
horizontal-source conclusions of that theorem, before any endpoint cap
or angular-turning estimate.  Its auxiliary growth cuts use the
qualified enlarged class and its constants from module 284.  The
stopped word classes below still use the original $k_0$ and its
original strips; these two partitions need not coincide.  Both
collections of source cuts will be avoided inside open physical
source intervals.  Lemma~\ref{lem:cs-determinant} gives
the exact two-half determinant inequality
\begin{equation}\label{ps:horizontal-product}
 u_x:=|\partial_{\alpha_0}p_N|^{-1}
 \le C_{\rm H}\cos\varphi_i\,a_I^-a_J^+,\qquad
 C_{\rm H}={\tau_-+R_+\over2\tau_-}={53\over12},
 \quad 1\le i<N,
\end{equation}
at a point of the actual overlap $I\cap J$.
Here the two half-lengths are $i$ and $N-i$.
At $i=0$ and $i=N$, respectively, the exact identities are
$u_x=|\partial_{\alpha_0}p_N|^{-1}$ on the forward horizontal source
and $u_x=|\partial_{\alpha_N}p_0|^{-1}$ on the backward horizontal
source.  The latter follows from preservation of
$d\alpha\,dp$, and is not a time-reversal substitution of a roof.

Corollary~\ref{cor:cs-overlap}, using the interval overlay lemma of
Section~\ref{sec:mb-capacity}, therefore gives at every common
horizontal source level
\begin{equation}\label{ps:horizontal-overlay}
 \sum_{O}\sup_Ou_x
 \le \cos h\,C_{\rm H}H_{\rm H}E_{\rm H}
                (\sigma^{i-1}+\sigma^{N-i-1}),
 \qquad 1\le i<N.
\end{equation}
For either endpoint the corresponding one-sided sum is at most
$H_{\rm H}$.  All source intervals fix actual lifted half-words.
Their overlaps consequently fix the complete physical word after
the initial disk is translated to the original fundamental copy.
They include arbitrarily short source intervals.

The marked-flight bound from Theorem~\ref{thm:mb-two-sided-sum}
is used only at positive clearance levels:
\begin{equation}\label{ps:marked-input}
 \sum_{O\in\mathcal O_{N,j,R,u}}\sup_Ou_x
 \le Q_{\rm mb}(\sigma^j+\sigma^{N-j-1}),
 \quad 0\le j<N,\quad 0<u\le3/200.
\end{equation}
Both mark signs and every actual source copy are already included
in this formula.

\subsection{An endpoint-exact comparison with the affine inverse weight}

Let $D^-_{N,R}$ be the open set of states with regular physical past
through times $0,-1,\ldots,-N$.  For $y\in D^-_{N,R}$ put
\begin{equation}\label{ps:affine-weight}
 e_R={ (1,-\kappa_R)\over\sqrt{1+\kappa_R^2}},
 \qquad
 w_N(y)={1\over |D T_R^{-N}(y)e_R|}.
\end{equation}
The vector is the unit tangent of the intrinsic affine fiber
$\Gamma_{R,\pi_R(y)}$ from Section~\ref{tc:section}.  The function
$w_N$ is positive and continuous on every finite regular word domain.
It does not require an infinite future or a choice of an invariant
stable line.

\begin{lemma}[The affine weight and the original endpoint momentum]
\label{ps:matrix}
For every regular physical $N$-flight record $x,T_Rx,\ldots,T_R^Nx$,
$N\ge1$, one has
\begin{equation}\label{ps:matrix-bound}
 w_N(T_R^Nx)\le R_+\Gamma\,c_0(x)\,
                          |\partial_{\alpha_0}p_N(x)|^{-1}.
\end{equation}
If $x\in M_R'$, then
\begin{equation}\label{ps:true-stable}
 j_N(x)\le C_\times w_N(T_R^Nx).
\end{equation}
Neither estimate has a terminal incidence cutoff.
\end{lemma}

\begin{proof}
Write the positive matrix product of Lemma~\ref{lem:eo-differential}
as
\[
 \mathsf D_N=\begin{pmatrix}A&B\\ C&D\end{pmatrix},\qquad
 D T_R^N=(-1)^N\mathsf D_N,\qquad
 \det\mathsf D_N=c_0/c_N.
\]
Every entry is positive.  Since a variation of $\alpha_0$ at fixed
$p_0$ has $\delta s_0=R\,\delta\alpha_0$ and
$\delta\varphi_0=0$, while $\delta p_N=c_N\delta\varphi_N$, the
original derivative is exactly
\begin{equation}\label{ps:u-matrix}
                 u_x={1\over R c_N C}.
\end{equation}
There is one radius factor and one terminal incidence in this
identity.

The inverse matrix sends $(1,-\kappa_R)$ to
\[
 {c_N\over c_0}
       (D+\kappa_R B,\,-C-\kappa_R A).
\]
Consequently
\begin{equation}\label{ps:w-matrix}
 w_N(T_R^Nx)
  ={c_0\sqrt{1+\kappa_R^2}\over
      c_N\sqrt{(D+\kappa_R B)^2+(C+\kappa_R A)^2}}.
\end{equation}
The square root in the denominator is at least $C$.
Dividing \eqref{ps:w-matrix} by $c_0u_x$ proves
\eqref{ps:matrix-bound}.  In particular no factor $1/c_N$ has
been dropped: the identical factor in \eqref{ps:u-matrix} cancels
before taking the uniform bound.

For \eqref{ps:true-stable}, true invariance gives
$j_N(x)=1/J_{E^s(T_R^Nx)}T_R^{-N}$.
At the same current point, both this stable direction and $e_R$
have slopes in $[-K,-a]$.  Apply
Corollary~\ref{cor:xd-backward-curves} with these two vectors at
the same point.  The physical word and all its past strip labels
agree automatically.  Inverting its two positive Jacobians gives
\eqref{ps:true-stable}.  No continuity assertion for an invariant
line near a singular point is involved.
\end{proof}

\subsection{The exhaustive physical exit of a fixed affine interval}

In this subsection a word is normalized by putting its current
scatterer at the origin.  Reversing its ordered centers and translating
its initial scatterer to the origin recovers one original forward
physical word, bijectively.  No periodic translate is counted twice.

Fix $N\ge1$.  Let $\mathcal U$ be the open subset of $D^-_{N,R}$
having a specified complete lifted past word and specified
\emph{open} homogeneity strips at exactly the times
$-1,\ldots,-N$.  There is no current homogeneity condition.
For any $y\in\mathcal U$ let $I_y$ be the connected open interval of
$\Gamma_{R,\pi_R(y)}\cap\mathcal U$ containing $y$.

\begin{lemma}[Physical boundary stopping]
\label{ps:exit}
At either endpoint of $I_y$, at least one of the following events
occurs in the limiting prescribed record, written in forward order
$0,\ldots,N$:
\begin{enumerate}
\item the current angle $\varphi_N$ is $+\pi/2$ or $-\pi/2$;
\item for some $0\le i<N$, $\varphi_i$ reaches a level
      $h\in\mathcal H$ bounding its fixed open strip;
\item a flight $0\le j<N$ is tangent in its interior to an unvisited
      lattice disk.
\end{enumerate}
All incidences $c_0,\ldots,c_{N-1}$ in this limiting record are
positive.  In the third alternative the tangent foot is strictly
inside the flight; no additional collision is inserted there.

For a chosen event, a tail of $I_y$ reaches arbitrarily small
positive interior values of its defining level distance while
remaining in the same complete word and all the same open strips.
Its image under that distance contains $(0,\eta)$ for some
$\eta>0$.  The number $\eta$ may depend on the word, the chosen
interval, and $R$.
\end{lemma}

\begin{proof}
Parametrize the intrinsic affine fiber by
$-\pi/2<\varphi_N<\pi/2$ as in \eqref{tc:fibers}.
An endpoint exists in the closed interval
$[-\pi/2,\pi/2]$, and the interval itself supplies an approaching
continuous path of actual regular records.  All listed centers are
fixed along it.  Their contact points lie on the corresponding
compact circles.  Every physical flight has length in
$[\tau_-,\tau_+]$.

Every one of the finitely many prescribed past strips has closure
strictly inside $(-\pi/2,\pi/2)$.  Thus
$c_0,\ldots,c_{N-1}$ have positive lower bounds along this interval;
these bounds are not asserted uniform in its strip labels.
Compactness gives a limiting prescribed record.  The limit is
unique: going backwards, the intersection with each prescribed
previous circle is a simple first-entry root because that circle's
limiting incidence is positive.  The implicit circle-intersection
formula consequently continues this prescribed branch uniquely
and continuously at each of the finitely many steps.  An unvisited
tangency does not obstruct this \emph{formal prescribed-word}
continuation.  No claim is made that the continued branch beyond
the endpoint remains a physical first-hit branch.

Limits of the actual closed flight segments avoid every open disk.
Only finitely many lattice disks can meet the uniformly bounded
neighborhood of any one of these finitely many segments.  Hence,
if the current angle is nongrazing, no past strip boundary is
reached, and no segment is tangent to an unvisited disk, every
listed collision remains a strict physical first collision on a
neighborhood of the endpoint.  All unvisited clearances there are
strict, and the finitely many prescribed strip conditions remain
open.  The endpoint would then belong to $\mathcal U$, and the
affine interval would extend through it.  This contradicts its
maximality and proves the exhaustive list.

An unvisited contact cannot be an endpoint of the flight, since
distinct radius-$R$ disks are disjoint with a positive gap.
A line segment avoiding an open disk can touch its boundary at an
interior point only tangentially.  This proves the stated nature
of the third event.  The geometric standoff of the original marked
source applies to this foot.

For a terminal grazing event take the interior distance
$\pi/2-\varepsilon\varphi_N$, with
$\varepsilon\in\{-1,1\}$ the selected terminal sign.
For a homogeneity event take
$\varepsilon(\varphi_i-h)$, where $\varepsilon$ is its side inside
the selected open strip.
For an unvisited tangency choose a sufficiently short tail on
which the foot remains inside the same flight and the signed
normal side of the marked disk is fixed; take the physical
clearance $u>0$ of this line from that disk.
In each case the distance is continuous and strictly positive on
the tail and tends to zero at the endpoint.  The continuous image
of an interval is an interval.  Since it contains one positive
value and values arbitrarily close to zero, it contains
$(0,\eta)$ for some $\eta>0$.
This is the intermediate-value theorem on actual records, with no
extension across the stopping boundary and no alteration of an
event test.
\end{proof}

\begin{lemma}[Common interior levels with actual source intervals]
\label{ps:common-level}
Select finitely many distinct physical words, one open past-strip
class and one stopped affine interval for each.  Assign each selected
word to one of its stopping events.  Group terminal events by their
sign; homogeneity events by $(i,h,\varepsilon)$; and unvisited
events by their marked flight index $j$, with both mark signs kept
on the original two source circles.

Within any such group, every sufficiently small common interior
level has a representative of each assigned word in its chosen
open class.  Each representative may be chosen away from all cuts
and exceptional sets of the one-sided source partitions.
Different assigned words then belong to different complete
overlay intervals, or different one-sided intervals in an
endpoint group.
\end{lemma}

\begin{proof}
For a finite group take the minimum of the positive level ranges
in Lemma~\ref{ps:exit}; for marked flights also keep the common
level at most $3/200$.  Its intermediate-value conclusion gives
the required physical representatives.  No uniform level
neighborhood over an infinite word family is used.

The relevant source slice is open in its one-dimensional source
parameter.  For an internal horizontal source this follows by
applying the regular physical map $T_R^i$ to an open neighborhood
of the selected record: the exact level $\varphi_i=h'$ cuts this
open set in an open $\alpha_i$ interval.  The same argument applies
at $i=0,N$.

For completeness, the marked source has the same property by an
explicit physical determinant.  At the departure state of flight
$j$, put $\vartheta=s_j/R+\varphi_j$, let $v$ be the outgoing
unit vector, and put
$b=(O-q_j)\cdot v^\perp$, $t=(O-q_j)\cdot v$ for the chosen
marked center $O$.  Direct differentiation gives
\[
 \partial_{s_j}\vartheta=\kappa_R,\quad
 \partial_{\varphi_j}\vartheta=1,\quad
 \partial_{s_j}b=-c_j-\kappa_R t,\quad
 \partial_{\varphi_j}b=-t,
 \qquad
 \det{\partial(\vartheta,b)\over\partial(s_j,\varphi_j)}=c_j>0.
\]
The signed clearance $u$ differs from $b$ by a fixed sign and the
constant $R$ on this tail.  Thus $(\vartheta,u)$ is a local
coordinate system there.  A fixed positive $u$ gives an open
angular interval inside the same class, with the same interior
foot and the same mark sign.

The excluded cuts and exceptional sets of either actual source
partition have one-dimensional Lebesgue measure zero.  Two such
null sets cannot exhaust a nonempty open source interval.
Choose the representative outside them.  It belongs to one
genuine one-sided interval in each required time direction.
Every such interval fixes its complete physical half-word;
an overlap fixes the full word.  The original one-copy
normalization follows the unique actual backward collisions to
the initial disk and translates that disk to the fundamental
copy.  Consequently two different assigned physical words
cannot occupy the same overlay interval.  The same argument with
one time direction proves the endpoint assertion.
\end{proof}

\subsection{The uniform unweighted sum}

Set
\begin{equation}\label{ps:source-budget}
 \mathcal B_*=
 2H_{\rm H}
 +2S_{\mathcal H}H_{\rm H}
 +{4S_{\mathcal H}C_{\rm H}H_{\rm H}E_{\rm H}\over1-\sigma}
 +{2Q_{\rm mb}\over1-\sigma},
 \qquad
 Q_{\rm stop}:=\max\{1,R_+\Gamma C_\times^2\mathcal B_*\}.
\end{equation}
All constants have been fixed by the physical geometry, original
homogeneity family, native source growth, and cross-curve theorem.
In particular $Q_{\rm stop}$ is independent of $N$, $R$, physical
word, exact record label, roof interval, and the outer parameters
of the raw local-limit problem.

\begin{theorem}[Full original physical-word pressure]
\label{ps:pressure}
For the original true stable Euclidean arclength cocycle,
\begin{equation}\label{ps:pressure-bound}
 P_N(R):=\sum_{w\in\mathcal W_N(R)}
          \sup_{x\in[w]_R\cap M_R'}j_N(x)
       \le Q_{\rm stop}
       \qquad(R\in[R_-,R_+],\ N\ge0).
\end{equation}
Every physical collision word is included.  No capacity abundance
premise, connectedness premise for a word domain, exclusion of
short cells, positive endpoint incidence cap, or pressure theorem
for a different model is used.
\end{theorem}

\begin{proof}
The case $N=0$ is the original convention $P_0=1$.
For $N\ge1$ let $\mathcal D_w$ be the finite regular past domain
of the current-normalized version of the forward word $w$.
Put
\[
 S_w=\sup_{y\in\mathcal D_w}w_N(y).
\]
A word with empty such domain has zero contribution.  Equation
\eqref{ps:true-stable} gives
\begin{equation}\label{ps:reduce-sup}
       P_N(R)\le C_\times\sum_w S_w.
\end{equation}
The finite cone expansion bounds these affine inverse Jacobians,
so each $S_w$ is finite.

Fix $0<\epsilon<1$ and any finite set of contributing physical
words.  For each choose $y_w$ off all past homogeneity boundaries
with
\[
                    w_N(y_w)\ge(1-\epsilon)S_w .
\]
This choice is legitimate: on the finite regular domain $w_N$
is continuous, and the pullback of every fixed homogeneity
level under a regular finite iterate is a smooth coordinate
level with empty interior.  The countably many such levels
have locally finite intersection with a neighborhood on which
the finitely many collision angles remain away from grazing.
Their complement is therefore dense.

Choose the open past-strip class containing $y_w$, stop its
affine interval by Lemma~\ref{ps:exit}, and assign one stopping
event.  For every actual representative $z$ in that same class,
the complete-word estimate of
Corollary~\ref{cor:xd-backward-curves} gives
\begin{equation}\label{ps:transfer}
 (1-\epsilon)S_w\le w_N(y_w)\le C_\times w_N(z)
                \le C_\times R_+\Gamma\,c_0(z)u_x(z).
\end{equation}
Only the matching past strips $-1,\ldots,-N$ are used here.
There is no current strip and no future invariant-line
regularity requirement on $z$.  In particular the transfer
does not follow a physical continuation from $y_w$ to $z$.

Sum \eqref{ps:transfer} by the finite source groups supplied by
Lemma~\ref{ps:common-level}.  It remains to bound the sum of
$c_0u_x$ of their selected representatives.

For a terminal grazing group, choose the common nongrazing
level $\varphi_N=\varepsilon(\pi/2-\delta)$.
The endpoint one-sided source bound, the uniqueness of the
assigned interval, and $c_0\le1$ pay the group by $H_{\rm H}$.
The two terminal signs cost $2H_{\rm H}$ in total.

For a homogeneity group $(i,h,\varepsilon)$ take the common
level $\varphi_i=h+\varepsilon\delta$.
If $i=0$, then $c_0=\cos(h+\varepsilon\delta)$ exactly, so the
one-sided bound pays its sum of $c_0u_x$ by
$\cos(h+\varepsilon\delta)H_{\rm H}$.
The selected $S_w$ on the left of \eqref{ps:transfer} do not
depend on $\delta$.  Letting this common positive parameter
decrease to zero gives the budget $\cos h\,H_{\rm H}$.
Both sides of every level and every level together cost
$2S_{\mathcal H}H_{\rm H}$.

For $1\le i<N$, use $c_0\le1$ and
\eqref{ps:horizontal-overlay}.  The same common-level argument
and the limit of its explicit cosine give the budget
\[
 \cos h\,C_{\rm H}H_{\rm H}E_{\rm H}
                       (\sigma^{i-1}+\sigma^{N-i-1})
\]
for that group.  The sum over the two sides, all levels, and all
internal indices is at most
\[
 2S_{\mathcal H}C_{\rm H}H_{\rm H}E_{\rm H}
       \sum_{i=1}^{N-1}(\sigma^{i-1}+\sigma^{N-i-1})
 \le {4S_{\mathcal H}C_{\rm H}H_{\rm H}E_{\rm H}\over1-\sigma}.
\]
This is where the actual collision incidence $\cos h$ is
essential: a uniform bound without it would not pay the
infinitely many original homogeneity boundaries.

For an unvisited group at index $j$, choose a common positive
clearance $u$ and use \eqref{ps:marked-input}.
The original two marked signs and the correct physical copy
are already included.  Its budget is
$Q_{\rm mb}(\sigma^j+\sigma^{N-j-1})$.  Thus all flight indices
cost at most
\[
 Q_{\rm mb}\sum_{j=0}^{N-1}(\sigma^j+\sigma^{N-j-1})
                      \le {2Q_{\rm mb}\over1-\sigma}.
\]
No zero-clearance trace, grazing derivative, or factor counting
unvisited disk labels is substituted into this step.

Combining the four positive budgets proves, for the chosen
finite physical word family,
\[
 \sum_w S_w\le
 {C_\times R_+\Gamma\over1-\epsilon}\mathcal B_*.
\]
Every selected word was assigned once.  Source partitions,
levels, and representatives may depend on this finite family,
but the displayed constant does not.  Nonnegative exhaustion
therefore proves the same inequality for the full word set;
in the present finite-horizon geometry there are already only
finitely many $N$-flight words.  Let $\epsilon\downarrow0$ and
apply \eqref{ps:reduce-sup}.  This gives
\eqref{ps:pressure-bound} with exactly \eqref{ps:source-budget}.
The common-level limits occurred at fixed $R,N$ and a fixed
finite selected family.  No physical collision-time limit,
roof essential supremum, or outer limit has been interchanged.
\end{proof}

\begin{corollary}[Completion of the original geometric pressure input]
\label{ps:input-complete}
Input~\ref{input:v112-geometric-pressure} holds on the original
radius family with
\[
        C_0=30000,\qquad C_1=6400,\qquad Q_*=Q_{\rm stop}.
\]
Consequently the formerly pressure-dependent constants in module 266
may be taken unconditionally as
\[
 C_\star={30000Q_{\rm stop}\over2c_*\lambda_\circ},
 \qquad K_\star=6400D Q_{\rm stop}.
\]
\end{corollary}
\begin{proof}
The two optical inequalities are
Theorem~\ref{thm:eo-optical}; the remaining full physical-word
sum is Theorem~\ref{ps:pressure}.
All parameters, measures, clock conventions, and exact-record
selectors in that input have been retained.
Thus Theorem~\ref{thm:v111-roof-pressure} and
Corollary~\ref{cor:v111-witness-moments} in module 266 apply with
these constants.  Their explicit descendants in modules 267--269,
namely Theorems~\ref{thm:v111-witness-payment},
\ref{thm:v112-load}, \ref{thm:v112-count}, and
\ref{thm:v112-payment}, likewise retain their stated conclusions
with $Q_*=Q_{\rm stop}$.  These are inherited consequences
activated by the newly proved geometric input; their proof bodies
and all their other hypotheses remain unchanged.
\end{proof}

\paragraph{Scope of the closure.}
This theorem pays the original unweighted physical-word pressure,
including the small-capacity word suprema left by
Section~\ref{tc:section}; it does not replace that capacity theorem.
The proof is independent of a planar-cell abundance hypothesis
and of a radius-uniform interpretation of a fixed-table external
pressure theorem.  The remaining spectral, arithmetic, conditional
microscopic-return, and raw pointwise local-limit obligations
retain their original quantifiers and order of limits.  A uniform
word-pressure bound is a completed geometric input to those
obligations, not a claim that their full conjunction has already
been proved.
In particular the normalized positive payments in modules
267--269 concern precisely the incidence events specified there.
They do not supply an $m^{-2}$ same-roof estimate for a single
unvisited-clearance witness, nor prove that the retained low-load
shortage in \eqref{eq:v112-remaining} vanishes in its original
ordered limits.

